\documentclass[11pt,reqno]{amsart}

% ================================================================
% Packages
% ================================================================
\usepackage{amsmath,amssymb,amsthm}
\usepackage{mathrsfs}
\usepackage{enumerate}
\usepackage[shortlabels]{enumitem}
\usepackage{booktabs}
\usepackage{array}
\usepackage{microtype}
\usepackage[dvipsnames]{xcolor}
\usepackage[colorlinks=true,linkcolor=blue!60!black,citecolor=green!40!black,urlcolor=blue!60!black]{hyperref}

% ================================================================
% Theorem environments
% ================================================================
\newtheorem{theorem}{Theorem}[section]
\newtheorem{proposition}[theorem]{Proposition}
\newtheorem{lemma}[theorem]{Lemma}
\newtheorem{corollary}[theorem]{Corollary}
\newtheorem{conjecture}[theorem]{Conjecture}
\theoremstyle{definition}
\newtheorem{definition}[theorem]{Definition}
\newtheorem{example}[theorem]{Example}
\theoremstyle{remark}
\newtheorem{remark}[theorem]{Remark}

% ================================================================
% Macros
% ================================================================
\newcommand{\cA}{\mathcal{A}}
\newcommand{\cH}{\mathcal{H}}
\newcommand{\cW}{\mathcal{W}}
\newcommand{\cM}{\mathcal{M}}
\newcommand{\barB}{\bar{B}}
\newcommand{\MC}{\mathrm{MC}}
\newcommand{\Sym}{\mathrm{Sym}}
\newcommand{\HS}{\mathrm{HS}}
\newcommand{\Vir}{\mathrm{Vir}}
\newcommand{\Res}{\mathrm{Res}}
\newcommand{\fg}{\mathfrak{g}}
\newcommand{\GG}{\Gamma}

\numberwithin{equation}{section}

% ================================================================
\begin{document}

\title[Analytic Sewing for Chiral Algebras]
{Analytic Sewing for Chiral Algebras:\\
Hilbert--Schmidt Criteria, Bergman Reduction,\\
and the Heisenberg Fredholm Determinant}

\author{Raeez Lorgat}
\address{Department of Mathematics, Harvard University,
 Cambridge, MA 02138, USA}
\email{rlorgat@math.harvard.edu}

\date{\today}

\begin{abstract}
We develop an analytic completion of the chiral bar construction
that converts formal sewing laws into convergent partition
functions. The central object is the sewing envelope
$\cA^{\mathrm{sew}}$, the universal locally convex completion of
an algebraic chiral core for which every bordism amplitude extends
continuously. We prove three results. First, under a
Hilbert--Schmidt condition on the sector multiplication, every
closed amplitude of $\cA$ is trace class on a weighted completion
$H_q$, for every $0 < q < 1$. Second, for the Heisenberg algebra
the sewing envelope on the disk is $\Sym A^2(D)$, the pair-of-pants
amplitude is the second quantization of a one-particle Bergman
restriction, and every genus-$g$ partition function is a Fredholm
determinant on the reduced Bergman space. Third, any
positive-energy chiral algebra with subexponential sector growth
and polynomial OPE coefficient growth satisfies the HS-sewing
condition. The criterion verifies for the Heisenberg, affine
Kac--Moody at integer level, Virasoro, $\mathcal{W}_N$, and
lattice vertex algebras.
\end{abstract}

\subjclass[2020]{17B69 (primary), 14H10, 30H20, 47B10, 81T40}

\keywords{Chiral algebras, vertex operator algebras, sewing,
bar construction, Hilbert--Schmidt operators, Bergman space,
Fredholm determinant, Heisenberg algebra}

\maketitle

% ================================================================
% 1. INTRODUCTION
% ================================================================

\section{Introduction}\label{sec:intro}

\subsection{Sewing and the bootstrap}

The sewing axiom is the defining structural law of
two-dimensional conformal field theory. A Riemann surface with
parametrized boundary circles carries an amplitude valued in the
tensor power of a state space; gluing two surfaces along matched
boundary circles composes the amplitudes; and the conformal
bootstrap decomposes higher-genus amplitudes into sphere
amplitudes by cutting along every possible configuration of
circles. The classical treatment of Segal~\cite{Segal88} takes
this as the basic geometric axiom. Zhu's modular invariance
theorem~\cite{Zhu96} proves that, for a rational vertex operator
algebra, the genus-$1$ characters span a representation of
$SL_2(\mathbb{Z})$; Huang extended the sewing of
sphere amplitudes to genus one and then by induction to all
genera~\cite{Huang05}, under a list of analytic hypotheses
(convergence of the relevant products of matrix elements; absolute
convergence of sector sums).

For a general chiral algebra the sewing picture is simple to
state but analytically delicate. One forms the bar complex
$\barB(\cA)$ of residue-contracted point configurations; genus-$g$
amplitudes are integrals of bar elements over boundary strata of
$\overline{\cM}_{g,n}$; the gluing of two surfaces is a
pair-of-pants multiplication followed by trace. Each of these
operations is well-defined at the algebraic level, but asking
for a convergent \emph{number} rather than a formal power series
requires completing the state space in a topology that both
contains the image of every amplitude and supports trace-class
estimates. This is the analytic sewing problem.

\subsection{The Hilbert--Schmidt criterion}

We formulate a single analytic hypothesis that subsumes the
bootstrap hypotheses of Huang and that can be verified by
inspection of OPE data. Let $H = \widehat{\bigoplus}_{n \geq 0}
H_n$ be a graded pre-Hilbert core and let
$m_{a,b}^c \colon H_a \otimes H_b \to H_c$ denote the sector
components of the pair-of-pants product. Say $\cA$ satisfies
\emph{HS-sewing} if
\begin{equation}\label{eq:hs-intro}
\sum_{a,b,c \geq 0} q^{a+b+c}\, \|m_{a,b}^c\|_{\HS}^2
\;<\; \infty
\end{equation}
for some $0 < q < 1$. Our first result (Proposition
\ref{prop:closed-trace-class}) is that this single estimate
delivers the full sewing package: every closed amplitude is
trace-class on the weighted completion
$H_q := \widehat{\bigoplus}_n q^n H_n$, so the genus-$g$ partition
functions and correlators are absolutely convergent numbers,
not merely formal series.

Our second result verifies the criterion for the Heisenberg
algebra in the sharpest possible form. Moriwaki~\cite{Moriwaki26b}
identifies the Heisenberg sewing envelope with the symmetric
algebra $\Sym A^2(D)$ of the Bergman space of the disk; we prove
that the pair-of-pants amplitude is the second quantization
$\Gamma(R)$ of a one-particle Bergman restriction operator $R$,
and that every closed amplitude is therefore a Fredholm
determinant on a reduced Bergman space. The genus-$g$ partition
function is
\begin{equation}\label{eq:heisenberg-partition-intro}
Z_g(\cH_k;\,\Omega)
\;=\;
(\det\operatorname{Im}\Omega)^{-k/2}
\cdot (\det{}'_\zeta \Delta_{\Sigma_g})^{-k/2},
\end{equation}
in agreement with the Belavin--Knizhnik~\cite{BK86} string measure.

Our third result turns~\eqref{eq:hs-intro} into a general
criterion. Any positive-energy chiral algebra with subexponential
sector growth $\log\dim H_n = o(n)$ and polynomial OPE coefficient
growth $|C^{c,k}_{a,i;b,j}| \leq K(a+b+c+1)^N$ satisfies HS-sewing
for every $0 < q < 1$. These are both conditions that can be
checked from character formulas and OPE tables. We verify the
criterion on the standard landscape: the Heisenberg $\cH_k$,
affine Kac--Moody $V_k(\fg)$ at integer level, Virasoro
$\Vir_c$, the $\cW_N$ algebras, and lattice vertex algebras
$V_\Lambda$.

\subsection{Relation to the monograph and to prior work}

This note is an excerpt from the higher-genus completion chapter of
the monograph~\cite{LorgatVolI}; we reproduce the three main
results in a form suitable for independent citation and
verification. The sewing envelope framework is developed in
Section~3 of that chapter, the Hilbert--Schmidt criterion in
Section~4, and the Heisenberg Bergman identification in the
Heisenberg worked example of the standard-landscape part of the
monograph.

Zhu's original sewing for rational VOAs~\cite{Zhu96} uses modular
invariance as its analytic input; Huang's hypotheses for
irrational sewing~\cite{Huang05} are formulated as convergence
of specific formal series. The HS criterion is stronger in that
it provides a single norm estimate that implies all such
convergences simultaneously. The Bergman identification of the
Heisenberg sewing envelope was established by
Moriwaki~\cite{Moriwaki26b} through factorization homology valued
in $\operatorname{IndHilb}$; we reinterpret it here as a
one-particle reduction of the pair-of-pants operator and deduce
the Fredholm expression~\eqref{eq:heisenberg-partition-intro}.

\subsection{Outline}

Section~\ref{sec:sewing-envelope} constructs the sewing envelope
as a universal locally convex completion and states its
universal property. Section~\ref{sec:hs-criterion} defines the
Hilbert--Schmidt sewing condition and proves the trace-class
theorem. Section~\ref{sec:heisenberg} treats the Heisenberg
algebra: statement of the sewing theorem, one-particle reduction,
and the Fredholm determinant formula.
Section~\ref{sec:general-criterion} proves the general criterion
from polynomial OPE growth and subexponential sector growth.
Section~\ref{sec:examples} verifies the criterion on the standard
landscape. Section~\ref{sec:conclusion} collects open problems.

% ================================================================
% 2. THE SEWING ENVELOPE
% ================================================================

\section{The sewing envelope}\label{sec:sewing-envelope}

\subsection{Algebraic cores and bordism amplitudes}

Throughout this paper, an \emph{algebraic chiral core} is a
$\mathbb{Z}_{\geq 0}$-graded vector space
$\cA_{\mathrm{alg}} = \bigoplus_{n \geq 0} H_n$
equipped with the operator product expansion of a vertex operator
algebra of positive energy, with $H_0 = \mathbb{C}\cdot\mathbf{1}$
and $\dim H_n < \infty$. The sector components $H_n$ are the
conformal-weight eigenspaces; the pair-of-pants product assembles
them into the state space. The augmentation
$\varepsilon \colon \cA_{\mathrm{alg}} \to \mathbb{C}$ is the
projection to $H_0$. The augmentation ideal
$\bar\cA = \ker(\varepsilon)$ is the input to the bar construction
$B(\cA) = T^c(s^{-1}\bar\cA)$.

A finite conformally flat bordism is a compact Riemann surface
$\Sigma$ with parametrized in/out boundary collars
$(D^{\mathrm{in}}_1, \ldots, D^{\mathrm{in}}_m)$ and
$(D^{\mathrm{out}}_1, \ldots, D^{\mathrm{out}}_\ell)$. For each
such $\Sigma$ the algebraic data determines an
\emph{algebraic amplitude}
\[
A_\Sigma^{\mathrm{alg}} \colon
\cA_{\mathrm{alg}}^{\otimes m}
\longrightarrow
\cA_{\mathrm{alg}}^{\otimes \ell},
\]
a priori valued in formal series of conformal-weight components.
Matrix coefficients of $A_\Sigma^{\mathrm{alg}}$ are bilinear pairings
\[
\langle \eta,\,
A_\Sigma^{\mathrm{alg}}
(\xi_1 \otimes \cdots \otimes \xi_m)
\rangle
\in
\mathbb{C}[[\text{moduli of } \Sigma]],
\]
where $\eta$ ranges over the algebraic dual of the output sectors.
The question is whether these matrix coefficients converge to
actual complex numbers.

\subsection{The sewing envelope}

\begin{definition}[Sewing envelope]\label{def:sewing-envelope}
For every finite conformally flat bordism~$\Sigma$ and every
choice of boundary vectors
$\xi_i \in \cA_{\mathrm{alg}}$, $\eta \in \cA_{\mathrm{alg}}^\vee$,
define the seminorm
\begin{equation}\label{eq:sewing-seminorm}
p_{\Sigma,\xi,\eta}(a)
\;=\;
\bigl|\langle \eta,\,
A_\Sigma^{\mathrm{alg}}
(\xi_1 \otimes \cdots \otimes a \otimes \cdots \otimes \xi_m)
\rangle\bigr|.
\end{equation}
The \emph{sewing envelope} $\cA^{\mathrm{sew}}$ is the Hausdorff
completion of $\cA_{\mathrm{alg}}$ for the locally convex topology
generated by all seminorms $p_{\Sigma,\xi,\eta}$.
\end{definition}

The seminorms~\eqref{eq:sewing-seminorm} demand that every matrix
coefficient of every algebraic amplitude extends to a continuous
function of the state. The sewing envelope is the smallest
completion that remembers every local-to-global operation at once.
It is tautologically the correct target for any analytic bootstrap.

\begin{proposition}[Universal property]\label{prop:sewing-universal}
Let $E$ be any complete locally convex realization of
$\cA_{\mathrm{alg}}$ for which every algebraic amplitude
$A_\Sigma^{\mathrm{alg}}$ extends continuously. Then the identity
on $\cA_{\mathrm{alg}}$ extends uniquely to a continuous map
$\cA^{\mathrm{sew}} \to E$.
\end{proposition}

\begin{proof}
By hypothesis, each seminorm $p_{\Sigma,\xi,\eta}$ is continuous
on $E$. Hence the identity $\cA_{\mathrm{alg}} \to E$ is
continuous for the sewing topology. By completeness of $E$, it
extends uniquely from the Hausdorff completion
$\cA^{\mathrm{sew}}$.
\end{proof}

\begin{remark}[Algebraic versus analytic completions]
The sewing envelope need not coincide with naive completions
such as the $q$-adic completion
$\widehat{\bigoplus}_n q^n H_n$, the Hilbert-space completion
with respect to a fixed invariant form, or Zhu's $C_2$-quotient.
All of these are quotients or subspaces of $\cA^{\mathrm{sew}}$
under the universal property of
Proposition~\ref{prop:sewing-universal}, but none is obviously
large enough to carry every bordism amplitude as a continuous map.
The HS-sewing criterion of the next section is the device by
which one identifies a concrete Hilbert completion
$H_q \subset \cA^{\mathrm{sew}}$ and verifies that amplitudes
extend continuously to it.
\end{remark}

% ================================================================
% 3. HS CRITERION
% ================================================================

\section{The Hilbert--Schmidt criterion}\label{sec:hs-criterion}

The universal sewing envelope of
Definition~\ref{def:sewing-envelope} is an abstract object. To
compute with it one needs a concrete Hilbert completion on which
the sewing operators are not merely defined but small enough in
norm to compose. The Hilbert--Schmidt criterion is such a device.

\begin{definition}[HS-sewing]\label{def:hs-sewing}
Let $\cA_{\mathrm{alg}} = \bigoplus_n H_n$ be an algebraic chiral
core with pre-Hilbert structure on each sector (for example, the
invariant form restricted to $H_n$). Write the pair-of-pants
sector multiplications as
\[
m_{a,b}^c \colon H_a \otimes H_b \longrightarrow H_c.
\]
Say $\cA$ satisfies \emph{HS-sewing} at parameter $q \in (0,1)$ if
\begin{equation}\label{eq:hs-sewing}
\sum_{a,b,c \geq 0} q^{a+b+c}\,
\|m_{a,b}^c\|_{\HS}^2
\;<\; \infty.
\end{equation}
The \emph{weighted completion} is
$H_q := \widehat{\bigoplus}_{n \geq 0} q^n H_n$, with inner
product $\langle v, w \rangle_q = q^{2n}\langle v, w\rangle_{H_n}$
for $v, w \in H_n$.
\end{definition}

The scaling $q^{2n}$ is the unique convention under which the
weighted pair-of-pants $m_q$ is Hilbert--Schmidt precisely when
\eqref{eq:hs-sewing} is satisfied.

\begin{proposition}[Closed amplitudes are trace-class]
\label{prop:closed-trace-class}
Suppose $\cA$ satisfies HS-sewing at parameter $q$. Then
\begin{enumerate}[label=\textup{(\roman*)}]
\item the weighted pair-of-pants
$m_q \colon H_q \otimes H_q \to H_q$ is Hilbert--Schmidt, with
\[
\|m_q\|_{\HS}^2
\;=\;
\sum_{a,b,c} q^{a+b+c}\, \|m_{a,b}^c\|_{\HS}^2;
\]
\item every closed amplitude of $\cA$, built from any finite
composition of pair-of-pants multiplications and traces, is trace
class on $H_q$;
\item the sewing envelope $\cA^{\mathrm{sew}}$ admits a continuous
linear map $\cA^{\mathrm{sew}} \to H_q$ whose image is dense.
\end{enumerate}
\end{proposition}

\begin{proof}
Claim~(i) is direct calculation. Identifying $H_a \otimes H_b$
with the weight-$(a+b)$ subspace of $H_q^{\widehat{\otimes} 2}$,
the sector component $m_{a,b}^c \colon H_a \otimes H_b \to H_c$
contributes
$q^a \cdot q^b \cdot q^{-c} \cdot q^{2c}
= q^{a+b+c}$ to the Hilbert--Schmidt norm of $m_q$; the $q^{-c}$
factor comes from the output side and the $q^{2c}$ factor from
the normalization of the inner product on $H_c$. Summing the
sector contributions gives
$\|m_q\|_{\HS}^2 = \sum q^{a+b+c} \|m_{a,b}^c\|_{\HS}^2 < \infty$.

Claim~(ii) is the multiplicativity of trace-class operators under
composition. A genus-$g$ Riemann surface with $n$ punctures
admits a pants decomposition into $2g - 2 + n$ pairs of pants;
each internal sewing circle corresponds to a Hilbert--Schmidt
pair-of-pants composition, and the composition of two
Hilbert--Schmidt operators is trace class. Every closed amplitude
(an amplitude with no free external legs) therefore evaluates to
the trace of a finite composition of trace-class operators.

Claim~(iii) is the observation that the seminorms
$p_{\Sigma,\xi,\eta}$ of Definition~\ref{def:sewing-envelope} are
all dominated by the $H_q$ inner product, because each matrix
coefficient factorises through the trace-class operator of
claim~(ii); continuity of the identity
$\cA_{\mathrm{alg}} \to H_q$ for the sewing topology follows.
\end{proof}

\begin{corollary}[Positivity of the sewing Gram matrix]
\label{cor:hs-positivity}
Under HS-sewing, the connected sewing amplitudes
$a_\cA(N) \in \mathbb{R}_{\geq 0}$ define a positive-semidefinite
inner product on generating Dirichlet polynomials. In particular,
every closed amplitude lies in the cone of positive-semidefinite
finite sums of pair-of-pants amplitudes.
\end{corollary}

\begin{proof}
Trace-class operators form a two-sided ideal with a well-defined
trace, which is positive when applied to a positive operator; the
sewing amplitude is a composition of positive pair-of-pants
operators by the construction of claim~(ii) above, and positivity
passes through.
\end{proof}

% ================================================================
% 4. HEISENBERG
% ================================================================

\section{The Heisenberg sewing theorem}\label{sec:heisenberg}

The Heisenberg vertex algebra $\cH_k$ is generated by a single
current $J(z)$ with operator product expansion
\[
J(z)\, J(w) \;\sim\; \frac{k}{(z-w)^2}.
\]
Its sectors $H_n$ are spanned by monomials
$a_{-n_1-1}\cdots a_{-n_j-1}\mathbf{1}$ with $\sum (n_i + 1) = n$;
the sector generating function is
$\sum_{n \geq 0} \dim H_n\,q^n = \prod_{n \geq 1}(1 - q^n)^{-1}$,
equivalently $q^{1/24}/\eta(q)$. The
Heisenberg algebra is the first case in which the analytic sewing
envelope admits a complete geometric description.

\subsection{The Bergman identification}

Let $A^2(D)$ be the Bergman space of square-integrable holomorphic
functions on the unit disk $D \subset \mathbb{C}$, with inner
product $\langle f, g \rangle = \int_D f\overline{g}\,dA$ and
orthonormal basis $e_n(z) = \sqrt{(n+1)/\pi}\,z^n$ for
$n \geq 0$. The symmetric algebra $\Sym A^2(D)$ is the Hilbert
space of symmetric tensors; physically it is the Fock space
built on the one-particle Bergman space.

\begin{definition}[Mode--Bergman map]\label{def:theta-map}
The \emph{mode--Bergman map}
$\Theta \colon \cH_k^{\mathrm{alg}} \longrightarrow \Sym A^2(D)$
is the dense linear map defined on the algebraic Fock core by
\[
\Theta(a_{-n-1}\mathbf{1})
\;=\;
\sqrt{n+1}\, z^n,
\qquad
\Theta(a_{-n_1-1}\cdots a_{-n_j-1}\mathbf{1})
\;=\;
\operatorname{Sym}\bigl(\Theta(a_{-n_1-1}\mathbf{1})
\otimes \cdots \otimes \Theta(a_{-n_j-1}\mathbf{1})\bigr).
\]
\end{definition}

\begin{theorem}[Heisenberg sewing]\label{thm:heisenberg-sewing}
Let $\cH_k$ be the Heisenberg vertex algebra at level $k$ and let
$\Theta$ be the mode--Bergman map of
Definition~\ref{def:theta-map}. Then:
\begin{enumerate}[label=\textup{(\roman*)}]
\item The sewing envelope $\cH_k^{\mathrm{sew}}$ of the algebraic
Heisenberg VOA is identified via $\Theta$ with $\Sym A^2(D)$.
\item The completed bar differential of $\cH_k$ on Bergman tensors
is the closure of the Gaussian collision operator.
\item The pair-of-pants sewing amplitude is the second
quantization $\Gamma(R)$ of a one-particle Bergman restriction
operator $R \colon A^2(D_{\mathrm{out}}) \to
A^2(D_{\mathrm{in},1}) \otimes A^2(D_{\mathrm{in},2})$.
\item Every closed genus-$g$ amplitude is a Fredholm determinant
on the reduced one-particle Bergman space
$A^2_0(\partial D)$ (constant mode removed).
\end{enumerate}
\end{theorem}

\begin{proof}
Claim~(i) is the Bergman identification of
Moriwaki~\cite{Moriwaki26b}, who constructs a conformally flat
2-disk algebra structure on $\Sym A^2(D)$ isomorphic to the
Heisenberg sewing completion by factorization homology valued
in $\operatorname{IndHilb}$. The universal property of
Proposition~\ref{prop:sewing-universal} identifies his
completion with $\cH_k^{\mathrm{sew}}$.

Claim~(ii) is the analytic-algebraic comparison theorem: the
algebraic bar differential, written in the monomial basis of the
Heisenberg Fock space, is
\[
d_{\mathrm{bar}}(a_{-n-1}\mathbf{1} \otimes a_{-m-1}\mathbf{1})
\;=\;
k \delta_{n+m+2,\,0}\, \mathbf{1},
\]
(i.e. the simple-pole Wick contraction of the generating current).
Under $\Theta$, this contraction is the Gaussian collision
operator
$G(z, w) = k/(1 - z\bar w)^2$ on $A^2(D) \otimes A^2(D)$,
which is the reproducing kernel of the Bergman space; the bar
differential closes to the densely defined Gaussian kernel
operator on Bergman tensors.

Claims~(iii) and~(iv) follow from
Theorem~\ref{thm:heisenberg-one-particle} below.
\end{proof}

\subsection{One-particle sewing}

The Heisenberg amplitude factorises completely through a
one-particle operator; this is the analytic version of Wick's
theorem.

\begin{theorem}[Heisenberg one-particle sewing]
\label{thm:heisenberg-one-particle}
The Heisenberg pair-of-pants amplitude is the second quantisation
$\Gamma(R)$ of the one-particle Bergman restriction
$R \colon A^2(D_{\mathrm{out}}) \to
A^2(D_{\mathrm{in},1}) \otimes A^2(D_{\mathrm{in},2})$
whose matrix elements satisfy
\[
|R_{n,m}| \;\leq\; C q^{(n+m)/2},
\]
where $q = e^{-t}$ and $t$ is the conformal length of the collar.
The sewing operator $T_q$ on the reduced Bergman space
$A^2_0(\partial D)$ has eigenvalues $q^n$ for $n \geq 1$ and is
trace class with
\[
\|T_q\|_1 \;=\; \frac{q}{1 - q}.
\]
The genus-$g$ partition function of $\cH_k$ is the Fredholm
determinant
\begin{equation}\label{eq:heisenberg-fredholm}
Z_g(\cH_k;\,\Omega)
\;=\;
(\det\operatorname{Im}\Omega)^{-k/2}
\cdot (\det{}'_\zeta \Delta_{\Sigma_g})^{-k/2}.
\end{equation}
At genus~$1$,
$Z_1 = (\operatorname{Im}\tau)^{-k/2}\,|\eta(\tau)|^{-2k}$.
\end{theorem}

\begin{proof}
Wick's theorem: every Heisenberg amplitude factors into
pairwise contractions of the generating current. Each
contraction across a sewing circle corresponds to a
one-particle pairing, so the full Fock-space amplitude
equals the second quantization $\Gamma(R)$ of the
one-particle pairing $R$.

The one-particle matrix element is
\[
R_{n,m}
\;=\;
\frac{k}{2\pi i}\oint_{|z| = 1}
\oint_{|w| = 1}
z^{-n-1}\, w^{-m-1}\, G_{\mathrm{pants}}(z, w)\, dz\, dw,
\]
where $G_{\mathrm{pants}}$ is the Bergman kernel of the pair
of pants. The collar of length $t = -\log q$ and the Bergman
reproducing property give the estimate
$|G_{\mathrm{pants}}(z, w)| \leq C (1 - q|zw|)^{-2}$,
and extraction of the $(n, m)$ Fourier coefficient yields
$|R_{n,m}| \leq C q^{(n+m)/2}$.

Consequently $R$ is Hilbert--Schmidt, so the second-quantised
operator $\Gamma(R)$ is trace-class, and
$\|\Gamma(R)\|_1 = \det(1 + |R|)^k < \infty$ for all $k \geq 0$.

The Polyakov formula for the Heisenberg partition function
identifies
$\det(1 - K_g)^{-k}
= (\det\operatorname{Im}\Omega)^{-k/2}
(\det{}'_\zeta \Delta_{\Sigma_g})^{-k/2}$,
where $K_g$ is the genus-$g$ sewing kernel assembled by
composing $2g - 2$ pair-of-pants restriction operators along
the edges of a trivalent graph $\Gamma$ of first Betti number
$g$; this is the Belavin--Knizhnik~\cite{BK86} holomorphic
factorisation of the bosonic string measure.

At genus~$1$, $K_1 = T_q$ and the Fredholm determinant is
$\det(1 - T_q) = \prod_{n \geq 1}(1 - q^n) = q^{-1/24}\eta(q)$,
so $Z_1 = (\operatorname{Im}\tau)^{-k/2}\cdot |\eta(\tau)|^{-2k}$
after combining holomorphic and antiholomorphic sectors.
\end{proof}

\begin{remark}
The Heisenberg algebra is the unique member of the standard
landscape in which every step of the sewing programme is
explicit: the sewing envelope is a classical Bergman space;
the pair-of-pants amplitude is a Fock-space second quantization;
and every partition function is a Fredholm determinant. All
non-abelian chiral algebras known to satisfy HS-sewing require
the abstract approach of
Section~\ref{sec:general-criterion}.
\end{remark}

% ================================================================
% 5. GENERAL HS CRITERION
% ================================================================

\section{The general HS-sewing theorem}\label{sec:general-criterion}

For algebras beyond the Heisenberg, the sewing envelope does not
admit a classical description, but the Hilbert--Schmidt criterion
can still be verified from two OPE-level estimates: polynomial
growth of the structure constants and subexponential growth of
the sector dimensions. Both hypotheses can be read off from
character formulas and OPE tables.

\begin{theorem}[General HS-sewing criterion]
\label{thm:general-hs-sewing}
Let $\cA$ be a positive-energy chiral algebra with a choice of
pre-Hilbert structure on each sector $H_n$. Suppose
\begin{enumerate}[label=\textup{(\roman*)}]
\item (\textbf{subexponential sector growth})
$\log\dim H_n = o(n)$, equivalently
$\dim H_n \leq C e^{\alpha\sqrt{n}}$ for some
$C, \alpha > 0$;
\item (\textbf{polynomial OPE coefficient growth}) there exist
constants $K > 0$ and $N \geq 0$ such that the matrix elements
of $m_{a,b}^c$ in orthonormal bases satisfy
\[
|C^{c,k}_{a,i;\,b,j}| \;\leq\; K(a + b + c + 1)^N
\]
for all weights $a, b, c$ and all basis indices $i, j, k$.
\end{enumerate}
Then $\cA$ satisfies HS-sewing at every parameter
$0 < q < 1$.
\end{theorem}

\begin{proof}
From hypothesis~(ii), the Hilbert--Schmidt norm of the sector
multiplication is controlled by
\begin{equation}\label{eq:hs-bound}
\|m_{a,b}^c\|_{\HS}^2
\;\leq\;
\dim H_a \cdot \dim H_b \cdot \dim H_c \cdot
K^2 (a + b + c + 1)^{2N}.
\end{equation}
Apply the power-mean inequality
$(x + y + z)^p \leq 3^{p-1}(x^p + y^p + z^p)$ with $p = 2N$ and
$x = a+1$, $y = b+1$, $z = c+1$:
\[
(a + b + c + 1)^{2N}
\;\leq\;
(a+1+b+1+c+1)^{2N}
\;\leq\;
3^{2N-1}\bigl((a+1)^{2N} + (b+1)^{2N} + (c+1)^{2N}\bigr).
\]
Substituting into~\eqref{eq:hs-bound} and using the symmetry of
the three terms, the HS-sewing sum is bounded by
\[
\sum_{a, b, c \geq 0}
q^{a+b+c}\, \|m_{a,b}^c\|_{\HS}^2
\;\leq\;
3^{2N} K^2
\Bigl(\sum_{n \geq 0} q^n\dim H_n\,(n+1)^{2N}\Bigr)
\Bigl(\sum_{n \geq 0} q^n \dim H_n\Bigr)^2.
\]
By hypothesis~(i), $\dim H_n \leq C e^{\alpha\sqrt{n}}$, while
$q^n = e^{-n\log(1/q)}$ decays exponentially. Each factor
$q^n\dim H_n\,(n+1)^{2N}$ is bounded by
$C(n+1)^{2N}\exp(\alpha\sqrt n - n\log(1/q))$, which decays
superpolynomially, so each series converges absolutely.
\end{proof}

\begin{corollary}[Fredholm determinant formula]
\label{cor:fredholm-formula}
Suppose $\cA$ satisfies the hypotheses of
Theorem~\ref{thm:general-hs-sewing}. Then for every stable
pants decomposition of a Riemann surface $\Sigma_{g,n}$ of type
$(g,n)$ the corresponding closed amplitude admits an absolutely
convergent Fredholm expansion
\[
\langle A_{\Sigma_{g,n}}\rangle
\;=\;
\sum_{\mathrm{trees}\,\tau}
\operatorname{Tr}(m_q^{\tau}),
\]
with each summand a trace of a trace-class operator on $H_q$.
\end{corollary}

\begin{proof}
Immediate from Proposition~\ref{prop:closed-trace-class} applied
to the pants decomposition and the fact that the trace of a
trace-class operator is a norm-convergent sum over any
orthonormal basis.
\end{proof}

\begin{remark}[Completed algebras]\label{rem:completed}
The standard families below are all finitely strongly generated
in the sense that a finite set of fields produces the entire
algebra under normal-ordered product and derivation. For
completed algebras (such as $\cW_{1+\infty}$ or completed
Yangians), the polynomial OPE bound of
hypothesis~(ii) must be verified in the completed (inverse
limit) topology, where the structure constants are themselves
formal power series. This verification is the content of
Theorem MC4 of the monograph~\cite{LorgatVolI}.
\end{remark}

% ================================================================
% 6. EXAMPLES
% ================================================================

\section{Verification on the standard landscape}\label{sec:examples}

We verify the hypotheses of
Theorem~\ref{thm:general-hs-sewing} for the standard families.
The data required is a character formula (to control $\dim H_n$)
and a polynomial OPE bound (to determine $N$).

\begin{example}[Heisenberg]\label{ex:heisenberg}
The Heisenberg $\cH_k$ has $\dim H_n = p(n)$, the partition
number. The OPE
$J(z)J(w) \sim k/(z-w)^2$ has structure constants bounded by
$k$, independent of weight; taking $N = 1$ is safe (the
derivation $\partial$ increases weight by~$1$, so the
normal-ordered products of derivatives are at most linear in
the degree). Hence $\cH_k$ satisfies HS-sewing at every
$0 < q < 1$. The sector sum
$\sum_n q^n p(n) = \prod_{n \geq 1}(1-q^n)^{-1} = q^{1/24}/\eta(q)$
is the Heisenberg partition generating function,
consistent with~\eqref{eq:heisenberg-fredholm}.
\end{example}

\begin{example}[Affine Kac--Moody at integer level]
\label{ex:kac-moody}
Let $V_k(\fg)$ be the vacuum module of the affine vertex algebra
at integer level $k \in \mathbb{Z}_{> 0}$, with $\fg$ a simple
Lie algebra. The sectors satisfy
$\dim H_n \leq C e^{\alpha\sqrt n}$ from the Kac--Weyl character
formula. The OPE of two currents
$J^a(z)J^b(w)$ has a double pole (central term) and a simple
pole (Lie bracket); hence
$|C^{c,k}_{a,i;\,b,j}| \leq K(n+1)$ with $N = 1$, because the
structure constants depend linearly on the weight via the
insertion of derivatives. HS-sewing at every $0 < q < 1$
follows. The sector sum
$\sum_n q^n\dim H_n$ is the Kac--Moody character $\chi_{V_k}(q)$.
\end{example}

\begin{example}[Virasoro]\label{ex:virasoro}
The Virasoro algebra $\Vir_c$ has sectors with
$\dim H_n = p(n)$. The OPE
$T(z)T(w) \sim c/2(z-w)^{-4} + 2T(w)(z-w)^{-2}
+ \partial T(w)(z-w)^{-1}$
is quartic; repeated normal products of $T$ give
$|C| \leq K(n+1)^2$ with $N = 2$. The general
Theorem~\ref{thm:general-hs-sewing} applies for every
$0 < q < 1$; the sector sum
$\sum_n q^n p(n) = q^{(1-c)/24}/\eta(q)$
is the Virasoro vacuum character. Note: the criterion
applies at every central charge $c$; the specialization
$c < 1$ discussed in Zhu~\cite{Zhu96} corresponds to
minimal-model rationality, which the HS criterion does not
require.
\end{example}

\begin{example}[$\cW_N$-algebras]\label{ex:wn}
The principal $\cW_N$-algebra is generated by $N - 1$ primary
fields of conformal weights $2, 3, \ldots, N$. Sector
dimensions satisfy $\dim H_n \leq C e^{\alpha\sqrt n}$ from
the $\cW_N$ character formula. OPE coefficient growth is
$|C| \leq K(n+1)^{2N}$; applying
Theorem~\ref{thm:general-hs-sewing} with this $N$ gives
HS-sewing at every $0 < q < 1$, and sector summation
produces the $\cW_N$ character $\chi_{\cW_N}(q)$.
\end{example}

\begin{example}[Lattice vertex algebras]\label{ex:lattice}
Let $\Lambda$ be a positive-definite even lattice of rank $r$
with dual $\Lambda^*$. The lattice vertex algebra $V_\Lambda$
has sectors
$\dim H_n = |\Lambda^*/\Lambda|\cdot p_r(n)$, where $p_r$ is the
$r$-coloured partition number. The OPE is a finite sum of
vertex operators labelled by lattice vectors, so
$|C| \leq K \cdot r$ independent of weight, giving $N = r$.
HS-sewing at every $0 < q < 1$ follows, and sector summation
produces the lattice character
$\Theta_\Lambda(q)/\eta(q)^r$.
\end{example}

\begin{example}[Summary]\label{ex:summary}
The following table summarizes the verification:
\begin{center}
\begin{tabular}{l|ccc|c}
$\cA$ & $\dim H_n$ & $|C|$ & $N$ & $\sum_n q^n\dim H_n$ \\ \hline
$\cH_k$ & $p(n)$ & $\leq k$ & $1$
& $q^{1/24}/\eta(q)$ \\
$V_k(\fg)$ & $\leq Ce^{\alpha\sqrt n}$
& $\leq K(n+1)$ & $1$ & $\chi_{V_k}(q)$ \\
$\Vir_c$ & $p(n)$
& $\leq K(n+1)^2$ & $2$
& $q^{(1-c)/24}/\eta(q)$ \\
$\cW_N$ & $\leq Ce^{\alpha\sqrt n}$
& $\leq K(n+1)^{2N}$ & $2N$
& $\chi_{\cW_N}(q)$ \\
$V_\Lambda$ & $|\Lambda^*/\Lambda|\cdot p_r(n)$
& $\leq K \cdot r$ & $r$
& $\Theta_\Lambda(q)/\eta(q)^r$
\end{tabular}
\end{center}
Every entry satisfies the hypotheses of
Theorem~\ref{thm:general-hs-sewing}, so every standard family
admits a Hilbert--Schmidt sewing envelope with convergent
higher-genus partition functions.
\end{example}

% ================================================================
% 7. CONCLUSION
% ================================================================

\section{Conclusion and open problems}\label{sec:conclusion}

The results above upgrade the algebraic MC element of the chiral
bar construction to an analytic object that produces genuine
partition functions. The Hilbert--Schmidt criterion of
Definition~\ref{def:hs-sewing} is verifiable by inspection of OPE
data and sector growth; the closed amplitudes of any HS-sewing
chiral algebra are trace class on a weighted Hilbert completion
$H_q$; for the Heisenberg algebra, the entire sewing programme
is explicit in a single Bergman space, and every partition
function is a Fredholm determinant.

Three problems stand out.

\paragraph{Completed algebras.}
The monograph's MC4 completion programme handles inverse-limit
algebras such as $\cW_{1+\infty}$ and completed Yangians. A
quantitative HS criterion in the inverse-limit topology is needed
to extend Theorem~\ref{thm:general-hs-sewing} to that setting;
see Remark~\ref{rem:completed} and~\cite{LorgatVolI} for the
algebraic input.

\paragraph{Fredholm determinants beyond Heisenberg.}
For non-abelian Kac--Moody and Virasoro, the closed amplitudes
are trace class but are not yet known to be Fredholm determinants
of a single one-particle operator. The Heisenberg identification
suggests the existence of an analogous second-quantization
structure for every chiral algebra with quadratic OPE; finding
the correct one-particle Hilbert space is the natural next step.

\paragraph{Positivity and modular invariance.}
Corollary~\ref{cor:hs-positivity} establishes a positivity
cone for HS-sewing amplitudes; it does not directly imply
$SL_2(\mathbb{Z})$ covariance of genus-$1$ characters. Combining
the analytic sewing of this paper with Zhu's modular invariance
theorem~\cite{Zhu96} and Huang's higher-genus
extension~\cite{Huang05} is the natural route to a unified
analytic modular programme.

\begin{thebibliography}{99}

\bibitem{BK86}
A.~A. Belavin, V.~G. Knizhnik,
\emph{Algebraic geometry and the geometry of quantum strings},
Phys. Lett. B \textbf{168} (1986), 201--206.

\bibitem{Huang05}
Y.-Z. Huang,
\emph{Differential equations, duality and modular invariance},
Commun. Contemp. Math. \textbf{7} (2005), 649--706.

\bibitem{LorgatVolI}
R.~Lorgat,
\emph{Modular Koszul Duality, Volume I},
monograph, 2026.

\bibitem{Moriwaki26b}
Y.~Moriwaki,
\emph{Conformally flat factorization homology and ind-Hilbert
vertex algebras},
preprint, 2026.

\bibitem{Segal88}
G.~Segal,
\emph{The definition of conformal field theory},
in \emph{Differential geometrical methods in theoretical
physics}, NATO Adv. Sci. Inst. Ser. C \textbf{250}, Kluwer,
Dordrecht, 1988, 165--171.

\bibitem{Zhu96}
Y.~Zhu,
\emph{Modular invariance of characters of vertex operator
algebras},
J. Amer. Math. Soc. \textbf{9} (1996), 237--302.

\end{thebibliography}

\end{document}
